% Gauss, Werke X/1 (1917), native134–135, printed127–128.
\SourceLeaf{134}\NoSourceHead
\pdfbookmark[1]{Allgemeines zur Lehre von den Gleichungen}{gleichungen-allgemeines}
\HistoricTitle{[ALLGEMEINES\\ZUR LEHRE VON DEN GLEICHUNGEN].}
\SourceGroup{V}
\HistoricTitle{[DIE BEZIEHUNGEN ZWISCHEN DEN KOEFFIZIENTEN\\
EINER GLEICHUNG UND DEN POTENZSUMMEN IHRER WURZELN.]}
\HistoricCitation{[Eintragung im \textsc{Leiste}, S. 6, 7.]}
\NumberedItem{1}
\noindent\hspace*{\parindent}Sit data aequatio
\[x^n-Ax^{n-1}+Bx^{n-2}-Cx^{n-3}+\ldots=0\]
\noindent cuius radices
\[x',\ x'',\ x''',\ldots\ x^N.\]
\noindent\hspace*{\parindent}Sit
\[\begin{aligned}
\Sigma x&=(1)\\
\Sigma xx&=(2)\\
\Sigma x^3&=(3)
\end{aligned}\]
\noindent\& sic deinceps: tum
\[\begin{aligned}
A&=(1)\\
B&=\tfrac12(1)^2-\tfrac12(2)\\
C&=\tfrac16(1)^3-\tfrac12(1)(2)+\tfrac13(3)\\
D&=\tfrac1{24}(1)^4-\tfrac14(1)^2(2)+\tfrac18(2)^2+\tfrac13(1)(3)-\tfrac14(4).
\end{aligned}\]
\noindent\hspace*{\parindent}Sit E. g.
\[(1)=-1;\quad(2)=-3;\quad(3)=+17;\quad(4)=-27.\]
\clearpage
\SourceLeaf{135}\ArithmeticHead{128}{L}{ALGEBRA.\quad NACHLASS.}
\noindent Erit aequatio fundam.
\[x^4+x^3+2xx-4x+3=0\]
\[(\text{Obiter }=(xx+\tfrac12x+\tfrac52)^2-13(\tfrac12x+\tfrac12)^2)\]
\NumberedItem{2}
\[\begin{aligned}
(1)&=A\\
(2)&=AA-2B\\
(3)&=A^3-3AB+3C\\
(4)&=A^4-4AAB+6B^2+4AC-4D\\
(5)&=A^5-5A^3B+5AB^2+5AAC-5BC-5AD+5E\\
(6)&=A^6-6A^4B+9A^2B^2-2B^3+6A^3C+3CC-12ABC\\
&\quad{}-6AAD+6AE-6F+6BD
\end{aligned}\]
\[\begin{aligned}
(n)&=A^n-nA^{n-2}B+\frac{n\mathpunct{.}n-3}{1\mathpunct{.}2}A^{n-4}B^2
-\frac{n\mathpunct{.}n-4\mathpunct{.}n-5}{1\mathpunct{.}2\mathpunct{.}3}A^{n-6}B^3\ldots\\[6pt]
&\ldots-nA^{n-3}C+\frac{n\mathpunct{.}n-5}{1\mathpunct{.}2}A^{n-6}CC
-\frac{n\mathpunct{.}n-7\mathpunct{.}n-8}{1\mathpunct{.}2\mathpunct{.}3}A^{n-9}C^3\ldots
\end{aligned}\]
\noindent legem generalem vide pagina sequenti.

Coefficiens termini $a^mb^nc^pd^q\ldots$ (si $A$ significet numerum permutationum
huius expressionis)
\[=\frac A{m+n+p+\ldots}(m+2n+3p+\ldots).\]
\begin{center}\textbf{BEMERKUNG.}\end{center}
\begingroup\small
\noindent\hspace*{\parindent}Zu der vorstehenden Aufzeichnung vergl. man die Tagebuchnotizen
Nr. 6 vom 23. Mai 1796 und Nr. 28 vom 21. August 1796 sowie die bezüglichen Bemerkungen
im Artikel 15 des \textsc{Bachmann}schen Aufsatzes »Über \textsc{Gauss}’ zahlentheoretische Arbeiten«.
\par\noindent\hfill\textsc{Schlesinger}.\par
\endgroup
\ClosingRule
\clearpage
